\section{Discussion and Limitations}
\label{sec:discussion}

\paragraph{Implications for hybrid design.}
For CNN-fed quantum heads the entangling layout and observables are choices about the backbone: product order, computable before training, ranks them (with a few reversals), and the best quantum heads are low-order.

\paragraph{Limitations.}
We study two families of CNOT layouts and a CZ ring (8 qubits, and 12 for the headline heads; $R_yR_z$ layers,
angle encodings, mainly $\pi\tanh$), one qubit per bottleneck coordinate, in noiseless simulation, with small CNNs and two pretrained networks on subsets
of MNIST, CIFAR-10 and SVHN of at most 10,000 images; the cost varies across pretrained backbones (1.2--7.2 points on CIFAR-10/4,000) and is untested with non-angle encodings
or on hardware. Most cells use 3 seeds (5 or 10 where stated), too few
for equivalence claims. Product order is a
coarse scalar: it misses brickwork by 5 points, its cost per unit does not carry across widths, and at fixed order the distribution of coupling also matters (\cref{tab:support}; classical heads only). The identity is exact only at zero rotations; trained small-CNN circuits move away from it (held-out $R^2$ 0.51--0.90). Both budgets assume independent factors and are loose at these orders; whether the smaller gradient mediates the cost, and how much of the intervention's gain is via saturation, are open. Information recoverable only by
decoders stronger than a two-hidden-layer MLP is not measured.
